% Werke X/1 (1917), native230–233, printed223–226.
\begingroup
\makeatletter\tagsleft@true\let\veqno\@@leqno\makeatother
\clearpage\SourceLeaf{230}\ArithmeticHead{223}{R}{\scriptsize DIE UNENDLICH VIELEN WERTE DES ARITHM.-GEOM.\ MITTELS.}
\par\indent Setzt man\par
\noindent\begin{minipage}[t]{.44\linewidth}
\[Pe^{-\pi t}=pt\]
\noindent so ist
\[\begin{aligned}
1)\quad&p(t+2ik)=pt\\[.7em]
2)\quad&pt=\frac1{\surd t}.p\frac1t
\end{aligned}\]
\end{minipage}\hfill\begin{minipage}[t]{.53\linewidth}
\[P(e^{-\pi t},\;e^{2\pi u})=p(t,u)\]
\noindent so ist
\[e^{-\frac{uu}{4\pi t}}\sqrt t.p(t,u)=p\left(\frac1t,\;\frac{iu}t\right)\]
\end{minipage}
\[\begin{gathered}
pt^4=p(t+i)^4+\left(\tfrac12p\tfrac14t-\tfrac12p(\tfrac14t+i)\right)^4\\[.6em]
pt=\tfrac12p\tfrac14t+\tfrac12p(\tfrac14t+i)\\[.6em]
pt+p(t+i)=2p4t\\[.6em]
p(t+i)=2p4t-pt=\frac1{\surd t}\left(p\frac1{4t}-p\frac1t\right).
\end{gathered}\]
\ClosingRule
\par\indent \(e^{-\pi}=h\) gesetzt, ist
\[\sum h^{\alpha nn+2\beta n+\gamma}=\frac{h^{\gamma-\frac{\beta\beta}\alpha}}{\surd\alpha}\sum h^{\frac1\alpha nn+\frac{2\beta i}\alpha.n}\]
\ClosingRule
\[\begin{array}{ll}
2ei+\frac1M=M'&pM'=\surd M.pM\\[.7em]
2di+\frac1{M'}=M''&pM''=\surd M'.pM'\\[.7em]
2ci+\frac1{M''}=M'''&pM'''=\surd M''.pM''\\[.7em]
2bi+\frac1{M'''}=M^{\mathrm{IV}}&\text{etc.}\\[.7em]
2ai+\frac1{M^{\mathrm{IV}}}=M^{\mathrm V}&\\[.4em]
\text{etc.}&
\end{array}\]
\noindent oder
\[\begin{aligned}
pM^{\mathrm V}&=\sqrt{MM'M''M'''M^{\mathrm{IV}}}.pM\\
pM^{\mathrm V}&=\surd[2bi,2ci,2di,2ei,M].pM.
\end{aligned}\]
\ClosingRule
\clearpage\SourceLeaf{231}\ArithmeticHead{224}{L}{ANALYSIS.\quad NACHLASS.}
\HistoricCitation{[Ein Blatt in Fg, Kapsel 46a.]}
\NumberedItem{7}
\par\indent Es sei
\[\alpha\delta-\beta\gamma=1\]
\noindent und
\[\frac{\alpha M-\beta i}{\delta+\gamma Mi}=N,\]
\[ac=bb+1,\quad AC=BB+1;\]
\noindent ist hier
\[M=\frac{1+bi}a,\quad N=\frac{1+Bi}A,\]
\noindent so geht die Form \((a,b,c)\) in \((A,B,C)\) über durch die Transf[ormation]
\[\begin{array}{rr}\delta&-\beta\\-\gamma&\alpha\end{array}\]
\par\indent Wir unterscheiden 6 Fälle jenachdem nach dem Modulus 2
\[\begin{array}{c|c|c|c|c|c|c}
\alpha&1&1&1&0&1&0\\
\beta&0&1&0&1&1&1\\
\gamma&0&0&1&1&1&1\\
\delta&1&1&1&1&0&0\\\hline
\multicolumn{1}{c}{}&\multicolumn{1}{c}{1}&\multicolumn{1}{c}{2}&\multicolumn{1}{c}{3}&\multicolumn{1}{c}{4}&\multicolumn{1}{c}{5}&\multicolumn{1}{c}{6}
\end{array}\]
\noindent Es ist dann
\begin{equation*}\tag*{[6]}
\begin{array}{@{}c@{}c|c|c|c|c|c|c}
&\multicolumn{1}{c}{}&\multicolumn{1}{c}{1}&\multicolumn{1}{c}{2}&\multicolumn{1}{c}{3}&\multicolumn{1}{c}{4}&\multicolumn{1}{c}{5}&\multicolumn{1}{c}{6}\\
\smash{\raisebox{-1\normalbaselineskip}{$\displaystyle\left\{\vphantom{\begin{array}{c}a\\a\\a\end{array}}\right.$}}&hpN=&pM&qM&rM&qM&rM&pM\\
&hqN=&qM&pM&pM&rM&pM&rM\\
&hrN=&rM&rM&qM&pM&qM&qM
\end{array}
\end{equation*}
\[h^{[-1]}=i^\alpha\surd(\delta+\gamma M).\]
\clearpage\SourceLeaf{232}\ArithmeticHead{225}{R}{\scriptsize DIE UNENDLICH VIELEN WERTE DES ARITHM.-GEOM.\ MITTELS.}
\NumberedItem{8}
\HistoricTitle{Die Reduction von \(pM,\;qM,\;rM\) auf die einfachste Form.}
\par\indent Es sei \(M=\dfrac{\alpha+\beta i}{\delta-\gamma i}\), wo \(\alpha,\beta,\gamma,\delta\) ganze reelle Zahlen und Zähler und Nenner ohne gemeinschaftlichen Factor; man setze
\[\begin{aligned}
\alpha\alpha+\beta\beta&=A\\
\alpha\gamma+\beta\delta&=B\\
\gamma\gamma+\delta\delta&=C\\
\alpha\delta-\beta\gamma&=\surd(AC-BB)=D.
\end{aligned}\]
\par\indent Man suche die einfachste Form des Det[erminanten] \(-DD\), welche der Form \((A,B,C)\) aequivalent ist; sie sei \((a,b,c)\).
\par\indent Dann lassen sich die Functionen von \(M\) auf Functionen von
\[\frac{D+bi}a\]
\noindent zurückführen. Der Algorithmus ist dieser
{\small\setlength{\arraycolsep}{8pt}
\[\begin{array}{lll}
\dfrac{D+Bi}A=M&DD+BB=AA'&B+B'=hA'\\[.9em]
\dfrac{D+B'i}{A'}=M'&DD+B'B'=A'A''&B'+B''=h'A''\\[.9em]
\dfrac{D+B''i}{A''}=M''&DD+B''B''=A''A''&B''+B'''=h''A'''
\end{array}\]}
\[\begin{array}{l}
\surd\dfrac{D+Bi}A.pM=pM'\text{ für gerades }h\\[.4em]
\hphantom{\surd\dfrac{D+Bi}A.pM}{}=qM'\text{ für ungerades }h\\[.5em]\hline\noalign{\vskip.5em}
\varepsilon^h\surd\dfrac{D+Bi}A.qM=rM'\qquad\varepsilon=\surd i\\[.5em]\hline\noalign{\vskip.5em}
\surd\dfrac{D+Bi}A.rM=qM'\text{ für gerades }h\\[.4em]
\hphantom{\surd\dfrac{D+Bi}A.rM}{}=pM'\text{ für ungerades }h
\end{array}\]
\par\indent Wenn man aus \(h,h',h''\), u.s.w.\ die Transformation von \((A,B,A')\) in \((a,b,c)\) ableitet, so werden deren Elemente (ob sie gerade oder ungerade) entscheiden, welche Funktion von \(\dfrac{D+bi}a\) mit der gegebenen von \(M\) so zusammenhängt, dass letztere in
\par\bigskip\noindent\small X1.\hfill29\normalsize
\clearpage\SourceLeaf{233}\ArithmeticHead{226}{L}{ANALYSIS.\quad NACHLASS.}
\[\varepsilon^H\surd\frac{D+Bi}A.\frac{D+B'i}{A'}.\frac{D+B''i}{A''}\ldots\]
\noindent multiplicirt werden muss. Wo \(M\) nicht rational ist, mag man \(D=-1\) setzen und den Algorithmus ebenso bilden. Nämlich wenn \(M=g+hi\), so geht man von der Form \(\dfrac1g,\;\dfrac hg,\;\dfrac{gg+hh}g\) (Det.\ \(-1\)) aus und sucht ihre Aequivalente etc.
\NumberedItem{9}
\par\indent Es ist
{\footnotesize\setlength{\arraycolsep}{4pt}\renewcommand{\arraystretch}{1.65}
\[\begin{array}{r|r|r|l}
\text{M. Ar.-G. zwischen }1\quad\ldots\;0{,}2&0{,}5208016&\log\ldots\;9{,}7166723&=\dfrac1{pM^2}\\
1{,}2\ldots\;0{,}8&0{,}9898721&\log\ldots\;9{,}9955791&=\dfrac1{\left(p\dfrac1M\right)^2}\\
\multicolumn{1}{r}{}&\multicolumn{1}{r}{0{,}5261302}&9{,}7210932&M
\end{array}\]}
\noindent Man sucht
{\small
\[\begin{aligned}
p(\tfrac17M+\tfrac67i)&=-0{,}4201578+0{,}3006159i\\[.4em]
&=q(\tfrac17M-\tfrac17i)\\[.4em]
&=r\left(\frac{7M}{1+MM}+\frac{7i}{1+MM}\right).\surd\frac{7(M+i)}{1+MM}
\end{aligned}\]
\[\begin{aligned}M&=\operatorname{cotg}\psi\\\psi&=62^0\,14'\,59''\\a&=\tfrac72\sin2\psi\\b&=7\sin\psi^2\end{aligned}\]
\[\begin{aligned}
r(a+bi)&=2A(\cos\tfrac14b\pi-i\sin\tfrac14b\pi)\\
&\quad+2A^9(\cos\tfrac94b\pi-i\sin\tfrac94b\pi)\\
&\qquad\text{etc.}
\end{aligned}\]}
\noindent\begin{minipage}[t]{.34\linewidth}
{\small\setlength{\arraycolsep}{2pt}
\[\begin{array}{r@{}r}
\multicolumn{2}{c}{\log A=-\frac14k\pi.a}\\[.4em]
\log\frac14k\pi={}&9{,}5328742\\
a\ldots{}&0{,}4600640\\\cline{2-2}
&9{,}9929382\\[.4em]
&0{,}9838711\\\cline{2-2}
A\ldots{}&9{,}0161289\\
2A\ldots{}&9{,}3171589\\
&0{,}39602\\\cline{2-2}
&9{,}71318\\
\multicolumn{2}{r}{\text{Richtig.}}
\end{array}\]}
\end{minipage}\hfill\vrule\hfill\begin{minipage}[t]{.61\linewidth}
\noindent Man hat
{\small
\[\int_0^\pi\frac{d\varphi}{\surd(1-\sin\varphi^2.\sin\lambda^2)}=\pi(pt)^2\]
\[\begin{aligned}
\int_0^\pi\frac{d\varphi}{\surd(1-\sin\varphi^2.\cos\lambda^2)}&=\pi\left(p\frac1t\right)^2\\
&=\pi t(pt)^2
\end{aligned}\]
\[\begin{aligned}qt&=pt.\surd\cos\lambda\\rt&=pt.\surd\sin\lambda.\end{aligned}\]}
\end{minipage}
\clearpage\endgroup
